\chapter{Matrix Multiplication and Inverses}

You already know what a single matrix \emph{does}: it takes a vector and moves it
somewhere. Feed in $\vec{x}$, get out $\mat{A}\vec{x}$. The matrix is a machine, and
the vector is the raw material that goes through it.

This chapter is about two questions that follow immediately.

First: what happens if you run a vector through \emph{two} machines in a row? Push
$\vec{x}$ through $\mat{B}$, then push the result through $\mat{A}$. You get
$\mat{A}(\mat{B}\vec{x})$. It turns out there is a single matrix --- we will call it
$\mat{A}\mat{B}$ --- that does the whole two-step journey in one shot. Finding that
matrix \emph{is} matrix multiplication. Everything about how you multiply matrices,
including the rule that feels so arbitrary the first time you see it, is forced on us
the moment we demand $(\mat{A}\mat{B})\vec{x} = \mat{A}(\mat{B}\vec{x})$.

Second: can we ever \emph{undo} a machine? If $\mat{A}$ sent $\vec{x}$ to $\vec{b}$,
is there a machine that takes $\vec{b}$ back to $\vec{x}$? That machine, when it exists,
is the \emph{inverse} $\mat{A}\inv$. We will see exactly when it exists, what it means
geometrically, and how to compute it with the same elimination you already know.

These two ideas are the grammar of everything that comes later. Factorizations like
$\mat{A}=\mat{L}\mat{U}$ and $\mat{A}=\mat{Q}\mat{R}$ are sentences in this grammar:
they say ``this complicated machine is really these simpler machines, one after another.''
The grammar is worth learning carefully.

\section{Where Matrix Multiplication Comes From}
\label{ch08:origin}

\subsection{Composition forces the rule}

Forget formulas for a moment. Here is the only thing we want.

\begin{intuitionbox}
A matrix is a machine that transforms vectors. If $\mat{B}$ is one machine and
$\mat{A}$ is another, then doing $\mat{B}$ \emph{first} and $\mat{A}$ \emph{second}
is itself some transformation. \textbf{Matrix multiplication is the name of that
combined transformation.} The product $\mat{A}\mat{B}$ is defined to be the single
matrix that produces $\mat{A}(\mat{B}\vec{x})$ for every $\vec{x}$.
\end{intuitionbox}

Notice the order already looks backwards: $\mat{B}$ acts first but is written on the
right, next to the $\vec{x}$ it touches first. That is not a typo and it is not a
convention someone could have chosen differently. It comes straight from how we write
function composition: $\,f(g(x))\,$ applies $g$ first, and $g$ sits on the inside,
closest to $x$. Matrices inherit that.

Let us make it concrete in two dimensions. Take

\[
\mat{B} = \begin{bmatrix} 2 & 0 \\ 0 & 1 \end{bmatrix}
\qquad\text{(stretch horizontally by }2\text{)},
\]
\[
\mat{A} = \begin{bmatrix} 0 & -1 \\ 1 & 0 \end{bmatrix}
\qquad\text{(rotate by }90^\circ\text{ counterclockwise)}.
\]

Take a test vector $\vec{x} = \colvec{1 \\ 1}$. Run it through $\mat{B}$ first:

\[
\mat{B}\vec{x} = \begin{bmatrix} 2 & 0 \\ 0 & 1 \end{bmatrix}\colvec{1 \\ 1}
= \colvec{2 \\ 1}.
\]

Now run \emph{that} through $\mat{A}$:

\[
\mat{A}(\mat{B}\vec{x}) = \begin{bmatrix} 0 & -1 \\ 1 & 0 \end{bmatrix}\colvec{2 \\ 1}
= \colvec{-1 \\ 2}.
\]

So the two-step journey sends $\colvec{1 \\ 1}$ to $\colvec{-1 \\ 2}$. Can we find a
\emph{single} matrix $\mat{C}$ that does this directly? Yes, and there is a quick
trick for reading it off. A matrix is completely determined by what it does to the
basis vectors $\vec{e}_1 = \colvec{1 \\ 0}$ and $\vec{e}_2 = \colvec{0 \\ 1}$, because
the columns of any matrix \emph{are} the images of the basis vectors. So just push
$\vec{e}_1$ and $\vec{e}_2$ through the whole pipeline and collect the answers as
columns.

\[
\mat{B}\vec{e}_1 = \colvec{2 \\ 0},\qquad
\mat{A}(\mat{B}\vec{e}_1) = \begin{bmatrix} 0 & -1 \\ 1 & 0 \end{bmatrix}\colvec{2\\0}
= \colvec{0 \\ 2}.
\]
\[
\mat{B}\vec{e}_2 = \colvec{0 \\ 1},\qquad
\mat{A}(\mat{B}\vec{e}_2) = \begin{bmatrix} 0 & -1 \\ 1 & 0 \end{bmatrix}\colvec{0\\1}
= \colvec{-1 \\ 0}.
\]

Stack those two images as columns:

\[
\mat{C} = \mat{A}\mat{B} = \begin{bmatrix} 0 & -1 \\ 2 & 0 \end{bmatrix}.
\]

Let us check it on our test vector: $\mat{C}\colvec{1\\1} = \colvec{0\cdot1 + (-1)\cdot1 \\ 2\cdot1 + 0\cdot1} = \colvec{-1 \\ 2}$. Exactly the two-step answer. The single
machine $\mat{C}$ reproduces ``stretch then rotate'' perfectly.

\begin{figure}[H]
\centering
\begin{tikzpicture}[scale=0.95,>=Stealth]
% three panels
\begin{scope}[xshift=0cm]
  \draw[->,gray!50] (-0.4,0)--(2.6,0);
  \draw[->,gray!50] (0,-0.4)--(0,2.6);
  \draw[->,thick,vecblue] (0,0)--(1,1) node[above right]{$\vec{x}$};
  \node at (1.2,-0.9) {start: $\colvec{1\\1}$};
\end{scope}
\begin{scope}[xshift=5cm]
  \draw[->,gray!50] (-0.4,0)--(2.6,0);
  \draw[->,gray!50] (0,-0.4)--(0,2.6);
  \draw[->,thick,structurepurple] (0,0)--(2,1) node[above right]{$\mat{B}\vec{x}$};
  \node at (1.2,-0.9) {after $\mat{B}$: $\colvec{2\\1}$};
\end{scope}
\begin{scope}[xshift=10cm]
  \draw[->,gray!50] (-1.4,0)--(1.6,0);
  \draw[->,gray!50] (0,-0.4)--(0,2.6);
  \draw[->,thick,vecred] (0,0)--(-1,2) node[above left]{$\mat{A}\mat{B}\vec{x}$};
  \node at (0.1,-0.9) {after $\mat{A}$: $\colvec{-1\\2}$};
\end{scope}
\draw[->,thick] (2.9,1.2)--(4.4,1.2) node[midway,above]{$\mat{B}$};
\draw[->,thick] (7.9,1.2)--(9.4,1.2) node[midway,above]{$\mat{A}$};
\end{tikzpicture}
\caption{Two machines in a row. First $\mat{B}$ stretches, then $\mat{A}$ rotates.
The single matrix $\mat{A}\mat{B}$ does the whole trip at once.}
\label{ch08:pipeline}
\end{figure}

\begin{focusbox}
To build $\mat{A}\mat{B}$, send each column of $\mat{B}$ through $\mat{A}$. The
$j$-th column of $\mat{A}\mat{B}$ is $\mat{A}$ times the $j$-th column of $\mat{B}$.
In symbols, if $\mat{B} = [\,\vec{b}_1 \ \vec{b}_2 \ \cdots \ \vec{b}_n\,]$, then
\[
\mat{A}\mat{B} = [\,\mat{A}\vec{b}_1 \ \ \mat{A}\vec{b}_2 \ \ \cdots \ \ \mat{A}\vec{b}_n\,].
\]
This is the single most useful way to think about a matrix product. Hold onto it.
\end{focusbox}

\subsection{Why the shapes have to match}

For $\mat{A}\mat{B}$ to make sense, the output of $\mat{B}$ must be a legal input for
$\mat{A}$. If $\mat{B}$ produces vectors in $\R^m$ (it has $m$ rows), then $\mat{A}$
must accept vectors in $\R^m$ (it has $m$ columns). That is the famous shape rule:

\[
\underbrace{\mat{A}}_{p \times m}\ \underbrace{\mat{B}}_{m \times n}
\;=\; \underbrace{\mat{A}\mat{B}}_{p \times n}.
\]

The inner dimensions ($m$ and $m$) must agree and they cancel; the outer dimensions
($p$ and $n$) survive as the shape of the product. There is nothing to memorize here
if you remember the machines: $\mat{B}$ maps $\R^n \to \R^m$, $\mat{A}$ maps
$\R^m \to \R^p$, so the composition maps $\R^n \to \R^p$.

\begin{figure}[H]
\centering
\begin{tikzpicture}[scale=0.9,font=\small,>=Stealth]
% A : p x m
\draw[fill=vecblue!12,draw=vecblue!70] (0,0) rectangle (1.6,2.4);
\node at (0.8,1.2) {$\mat{A}$};
\node[below] at (0.8,-0.05) {$m$};
\node[left] at (-0.05,1.2) {$p$};
% B : m x n
\draw[fill=structurepurple!14,draw=structurepurple!80] (2.4,0.4) rectangle (4.4,2.0);
\node at (3.4,1.2) {$\mat{B}$};
\node[below] at (3.4,0.35) {$n$};
\node[left] at (2.35,1.2) {$m$};
\node at (5.2,1.2) {$=$};
% AB : p x n
\draw[fill=vecred!12,draw=vecred!70] (6.0,0) rectangle (8.0,2.4);
\node at (7.0,1.2) {$\mat{A}\mat{B}$};
\node[below] at (7.0,-0.05) {$n$};
\node[left] at (5.95,1.2) {$p$};
% inner-dimension match callout
\draw[->,thick,targetgreen] (1.7,-0.65) .. controls (1.9,-1.05) and (2.3,-1.05) .. (2.5,-0.65);
\node[targetgreen,below] at (2.1,-1.0) {inner $m$ must match, then cancels};
% outer dimensions survive
\draw[->,gray!70,dashed] (0.4,2.7) -- (6.4,2.7);
\node[gray!70,above] at (3.4,2.7) {outer $p$ and $n$ survive};
\end{tikzpicture}
\caption{The shape rule as a picture. The right edge of $\mat{A}$ ($m$ columns) must
meet the left edge of $\mat{B}$ ($m$ rows); the shared inner dimension cancels, and the
outer dimensions $p$ and $n$ become the shape of $\mat{A}\mat{B}$.}
\label{ch08:shaperule}
\end{figure}

\begin{trapbox}
``$\mat{A}\mat{B}$ exists'' does \emph{not} mean ``$\mat{B}\mat{A}$ exists.'' If
$\mat{A}$ is $2\times 3$ and $\mat{B}$ is $3\times 4$, then $\mat{A}\mat{B}$ is
$2\times 4$ --- fine --- but $\mat{B}\mat{A}$ asks a $3\times4$ to feed a $2\times3$,
and $4 \ne 2$, so it is undefined. Order matters before we even get to the values.
\end{trapbox}

\section{Four Views of the Same Product}
\label{ch08:fourviews}

There is one product $\mat{A}\mat{B}$, but four ways to see it, and a fluent reader
can switch between them at will. Each view illuminates a different chapter later in the
book, so we walk through all four with the same small example. Throughout this section
let

\[
\mat{A} = \begin{bmatrix} 1 & 2 \\ 3 & 4 \end{bmatrix},
\qquad
\mat{B} = \begin{bmatrix} 5 & 6 \\ 7 & 8 \end{bmatrix}.
\]

We will compute $\mat{A}\mat{B}$ four times and get the same answer each time:

\[
\mat{A}\mat{B} = \begin{bmatrix} 19 & 22 \\ 43 & 50 \end{bmatrix}.
\]

\subsection{View 1: the dot-product (entry) view}

This is the rule you will use to grind out numbers by hand. The entry in row $i$,
column $j$ of $\mat{A}\mat{B}$ is the dot product of row $i$ of $\mat{A}$ with column
$j$ of $\mat{B}$:

\[
(\mat{A}\mat{B})_{ij} = \sum_{k} \mat{A}_{ik}\,\mat{B}_{kj}
= (\text{row } i \text{ of } \mat{A}) \cdot (\text{column } j \text{ of } \mat{B}).
\]

Let us do all four entries of our example, slowly.

\[
(\mat{A}\mat{B})_{11} = \colvec{1 \\ 2}\cdot\colvec{5 \\ 7} = 1\cdot5 + 2\cdot7 = 5+14 = 19,
\]
\[
(\mat{A}\mat{B})_{12} = \colvec{1 \\ 2}\cdot\colvec{6 \\ 8} = 1\cdot6 + 2\cdot8 = 6+16 = 22,
\]
\[
(\mat{A}\mat{B})_{21} = \colvec{3 \\ 4}\cdot\colvec{5 \\ 7} = 3\cdot5 + 4\cdot7 = 15+28 = 43,
\]
\[
(\mat{A}\mat{B})_{22} = \colvec{3 \\ 4}\cdot\colvec{6 \\ 8} = 3\cdot6 + 4\cdot8 = 18+32 = 50.
\]

\begin{intuitionbox}
``Row $i$ dotted with column $j$.'' To find a single entry you reach across the row of
the left matrix and down the column of the right matrix, multiply pairwise, and add.
It is the most mechanical view, perfect for hand computation, but the worst for
\emph{understanding} --- it hides the geometry. Use the next three views for meaning.
\end{intuitionbox}

\begin{marginnotebox}
\textbf{The cost.} This view also lets us count the work. Multiplying a $p\times m$ by an
$m\times n$ matrix fills $p\,n$ entries, and each entry is a dot product of length $m$:
about $m$ multiplications and $m$ additions. So a full product costs on the order of
$p\,m\,n$ multiply--adds, and for two $n\times n$ matrices that is roughly $n^3$
operations. Keep that $n^3$ in mind; it is why the order you multiply a long chain
matters, and why the rank-one and factored views below are not just pretty but
\emph{cheaper}.
\end{marginnotebox}

\subsection{View 2: the column view (combinations of $\mat{A}$'s columns)}

This is the view from Section~\ref{ch08:origin}: each column of $\mat{A}\mat{B}$ is
$\mat{A}$ acting on the corresponding column of $\mat{B}$. And remember what $\mat{A}$
times a vector means --- it is a \emph{linear combination of the columns of $\mat{A}$},
with weights given by the entries of that vector.

Write $\mat{A} = [\,\vec{a}_1 \ \vec{a}_2\,]$ with $\vec{a}_1 = \colvec{1\\3}$,
$\vec{a}_2 = \colvec{2\\4}$. Column $1$ of $\mat{B}$ is $\colvec{5\\7}$, so:

\[
\text{col 1 of } \mat{A}\mat{B} \;=\; 5\,\vec{a}_1 + 7\,\vec{a}_2
= 5\colvec{1\\3} + 7\colvec{2\\4}
= \colvec{5\\15} + \colvec{14\\28} = \colvec{19\\43}.
\]

And column $2$ of $\mat{B}$ is $\colvec{6\\8}$, so:

\[
\text{col 2 of } \mat{A}\mat{B} \;=\; 6\,\vec{a}_1 + 8\,\vec{a}_2
= 6\colvec{1\\3} + 8\colvec{2\\4}
= \colvec{6\\18} + \colvec{16\\32} = \colvec{22\\50}.
\]

Same matrix $\begin{bmatrix}19 & 22\\ 43 & 50\end{bmatrix}$, built column by column.

\begin{focusbox}
\textbf{Every column of $\mat{A}\mat{B}$ lives in the column space of $\mat{A}$.}
Each one is a combination of $\mat{A}$'s columns. This single fact is why
$\col(\mat{A}\mat{B}) \subseteq \col(\mat{A})$ and therefore
$\rank(\mat{A}\mat{B}) \le \rank(\mat{A})$ --- you cannot create new directions by
multiplying on the right. We will lean on this hard when we study rank and the four
subspaces.
\end{focusbox}

\subsection{View 3: the row view (combinations of $\mat{B}$'s rows)}

By perfect symmetry, each \emph{row} of $\mat{A}\mat{B}$ is a combination of the
\emph{rows} of $\mat{B}$, with weights from the corresponding row of $\mat{A}$.

Row $1$ of $\mat{A}$ is $[\,1 \ \ 2\,]$. Rows of $\mat{B}$ are
$\vec{r}_1 = [\,5 \ \ 6\,]$ and $\vec{r}_2 = [\,7 \ \ 8\,]$:

\[
\text{row 1 of } \mat{A}\mat{B} = 1\cdot[\,5 \ 6\,] + 2\cdot[\,7\ 8\,]
= [\,5\ 6\,] + [\,14\ 16\,] = [\,19\ \ 22\,].
\]

Row $2$ of $\mat{A}$ is $[\,3\ \ 4\,]$:

\[
\text{row 2 of } \mat{A}\mat{B} = 3\cdot[\,5\ 6\,] + 4\cdot[\,7\ 8\,]
= [\,15\ 18\,] + [\,28\ 32\,] = [\,43\ \ 50\,].
\]

\begin{focusbox}
Mirror image of View 2: \textbf{every row of $\mat{A}\mat{B}$ lives in the row space
of $\mat{B}$.} Multiplying on the \emph{left} by $\mat{A}$ cannot create new row
directions, so $\rank(\mat{A}\mat{B}) \le \rank(\mat{B})$ too. Combine with View 2 and
you get $\rank(\mat{A}\mat{B}) \le \min(\rank \mat{A}, \rank \mat{B})$ for free.
\end{focusbox}

\subsection{View 4: the outer-product (sum-of-rank-one) view}

Here is the view that looks strangest and pays off the most later. Instead of building
the answer entry by entry, build it as a \emph{sum of whole matrices}. Multiply column
$k$ of $\mat{A}$ by row $k$ of $\mat{B}$ --- a column times a row is a full matrix, an
\emph{outer product} --- and add up over $k$:

\[
\mat{A}\mat{B} = \sum_{k} (\text{col } k \text{ of } \mat{A})(\text{row } k \text{ of } \mat{B}).
\]

A column times a row, like $\colvec{1\\3}[\,5 \ 6\,]$, fills out a whole matrix:

\[
\colvec{1\\3}\,[\,5\ \ 6\,] = \begin{bmatrix} 1\cdot5 & 1\cdot6 \\ 3\cdot5 & 3\cdot6 \end{bmatrix}
= \begin{bmatrix} 5 & 6 \\ 15 & 18 \end{bmatrix}.
\]

That is an outer product. It has rank one: every row is a multiple of $[\,5\ 6\,]$,
and every column is a multiple of $\colvec{1\\3}$. Now form the second one from column
$2$ of $\mat{A}$ and row $2$ of $\mat{B}$:

\[
\colvec{2\\4}\,[\,7\ \ 8\,] = \begin{bmatrix} 14 & 16 \\ 28 & 32 \end{bmatrix}.
\]

Add the two rank-one pieces:

\[
\mat{A}\mat{B} = \begin{bmatrix} 5 & 6 \\ 15 & 18 \end{bmatrix}
+ \begin{bmatrix} 14 & 16 \\ 28 & 32 \end{bmatrix}
= \begin{bmatrix} 19 & 22 \\ 43 & 50 \end{bmatrix}.
\]

Same answer, fourth time. The product is a \emph{layered sum of simple rank-one slabs}.

\begin{figure}[H]
\centering
\begin{tikzpicture}[scale=0.62,font=\small]
\tikzset{cell/.style={draw=gray!55,minimum size=0.95cm}}
% outer product 1
\matrix(P1)[matrix of nodes,nodes={cell},column sep=-\pgflinewidth,row sep=-\pgflinewidth] at (0,0)
{ 5 & 6 \\ 15 & 18 \\};
\node[above=2pt of P1]{$\vec{a}_1\,\vec{r}_1$ (rank 1)};
\node at (3.0,-0.45){$+$};
% outer product 2
\matrix(P2)[matrix of nodes,nodes={cell},column sep=-\pgflinewidth,row sep=-\pgflinewidth] at (4.8,0)
{ 14 & 16 \\ 28 & 32 \\};
\node[above=2pt of P2]{$\vec{a}_2\,\vec{r}_2$ (rank 1)};
\node at (7.8,-0.45){$=$};
% result
\matrix(R)[matrix of nodes,nodes={cell},column sep=-\pgflinewidth,row sep=-\pgflinewidth] at (9.6,0)
{ 19 & 22 \\ 43 & 50 \\};
\node[above=2pt of R]{$\mat{A}\mat{B}$};
\end{tikzpicture}
\caption{The outer-product view: $\mat{A}\mat{B}$ is a sum of rank-one layers, one per
column of $\mat{A}$ paired with the matching row of $\mat{B}$.}
\label{ch08:outerprod}
\end{figure}

\begin{applicationbox}
\textbf{Low-rank approximation and recommender systems.} The outer-product view is the
gateway to the most important idea in modern data science: a big matrix can often be
\emph{well approximated} by a short sum of rank-one layers. The singular value
decomposition (Chapter on the SVD) writes any matrix as
$\mat{A} = \sum_k \sigma_k\, \vec{u}_k \vec{v}_k\T$, a sum of outer products weighted by
the singular values $\sigma_k$. Keeping only the few largest terms gives the best
low-rank approximation. When Netflix predicts how you would rate an unseen film, it is
modeling the giant users-by-movies rating matrix as a short sum of such layers: a few
hidden ``taste directions'' $\vec{u}_k$ times a few ``genre directions'' $\vec{v}_k\T$.
Image compression, latent semantic analysis in search, and the embeddings inside
language models all live in this same picture. View 4 is where it starts.
\end{applicationbox}

\begin{walkbox}
Take a walk and try to feel why all four views must agree. Here is a hint to chew on:
the entries $(\mat{A}\mat{B})_{ij} = \sum_k \mat{A}_{ik}\mat{B}_{kj}$ share the same
sum over $k$. If you fix $j$ and collect all $i$, you get columns (View 2). If you fix
$i$ and collect all $j$, you get rows (View 3). And if you pull out one $k$ at a time across all $i,j$ at
once, you get the rank-one layers (View 4). Four views, one shared sum, freedom to
group it however you like.
\end{walkbox}

\section{The Algebra of Multiplication}
\label{ch08:algebra}

Now that we know what the product \emph{is}, what rules does it obey? Some are
comforting, one is famously not.

\subsection{It is associative (and the reason is composition)}

\begin{theorem}[Associativity]
\label{ch08:assoc}
For any matrices of compatible shapes,
\[
(\mat{A}\mat{B})\mat{C} = \mat{A}(\mat{B}\mat{C}).
\]
\end{theorem}

You could prove this by writing out triple sums of subscripts, and it works, but it is
joyless. The real reason is composition. Each matrix is a function. $\mat{A}\mat{B}\mat{C}$
means ``do $\mat{C}$, then $\mat{B}$, then $\mat{A}$,'' and \emph{that ordered sequence
of three machines is the same sequence no matter how you mentally bracket it.} Grouping
the first two or the last two does not change which machine runs when. Associativity of
matrix multiplication is just associativity of ``and then'' for functions.

\begin{formalbox}
\textbf{Why composition is associative.} For functions, $((f\circ g)\circ h)(x) =
(f\circ g)(h(x)) = f(g(h(x)))$ and $(f\circ(g\circ h))(x) = f((g\circ h)(x)) =
f(g(h(x)))$. Both equal $f(g(h(x)))$. Since a matrix is exactly the function
$\vec{x}\mapsto \mat{A}\vec{x}$, and matrix product is composition, associativity is
inherited. No subscripts required.
\end{formalbox}

\begin{aestheticbox}
Associativity quietly does enormous work. It is the reason a chain like $\mat{A}\mat{B}\mat{C}\mat{D}$ is unambiguous, so factorizations ($\mat{L}\mat{U}$, $\mat{Q}\mat{R}$,
$\mat{U}\mat{\Sigma}\mat{V}\T$) are even meaningful. It is the reason you can choose the
\emph{cheapest order} to multiply a long chain --- multiplying $\mat{A}(\mat{B}\vec{x})$
as matrix-times-vector costs far less than forming $\mat{A}\mat{B}$ first --- a trick
that powers efficient deep-learning forward passes. One small algebraic law, paid out
across the entire subject.
\end{aestheticbox}

\subsection{It is NOT commutative}

\begin{theorem}[Order matters]
In general $\mat{A}\mat{B} \ne \mat{B}\mat{A}$.
\end{theorem}

The cleanest proof is one example, and our opening machines already do it. Rotating then
stretching is genuinely different from stretching then rotating. Let

\[
\mat{A} = \begin{bmatrix} 0 & -1 \\ 1 & 0 \end{bmatrix}\ (\text{rotate } 90^\circ),
\qquad
\mat{B} = \begin{bmatrix} 2 & 0 \\ 0 & 1 \end{bmatrix}\ (\text{stretch } x).
\]

We already found $\mat{A}\mat{B} = \begin{bmatrix} 0 & -1 \\ 2 & 0 \end{bmatrix}$.
Now the other order:

\[
\mat{B}\mat{A} = \begin{bmatrix} 2 & 0 \\ 0 & 1 \end{bmatrix}
\begin{bmatrix} 0 & -1 \\ 1 & 0 \end{bmatrix}
= \begin{bmatrix} 0 & -2 \\ 1 & 0 \end{bmatrix}.
\]

These are different matrices. Watch them act on $\vec{e}_1 = \colvec{1\\0}$. With
$\mat{A}\mat{B}$: stretch $\vec{e}_1$ to $\colvec{2\\0}$, then rotate to $\colvec{0\\2}$.
With $\mat{B}\mat{A}$: rotate $\vec{e}_1$ to $\colvec{0\\1}$, then stretch the
\emph{x}-coordinate (which is already $0$) to get $\colvec{0\\1}$. One ends at
$\colvec{0\\2}$, the other at $\colvec{0\\1}$. The stretch lands on a different
direction depending on whether the rotation happened before or after it.

\begin{figure}[H]
\centering
\begin{tikzpicture}[scale=1.0,>=Stealth]
\begin{scope}[xshift=0cm]
  \draw[->,gray!50] (-0.4,0)--(2.6,0);
  \draw[->,gray!50] (0,-0.4)--(0,2.6);
  \draw[->,thick,vecblue] (0,0)--(0,2) node[right]{$\colvec{0\\2}$};
  \node[align=center] at (1.3,-1.0) {$\mat{A}\mat{B}\,\vec{e}_1$\\ stretch \emph{then} rotate};
\end{scope}
\begin{scope}[xshift=6cm]
  \draw[->,gray!50] (-0.4,0)--(2.6,0);
  \draw[->,gray!50] (0,-0.4)--(0,2.6);
  \draw[->,thick,vecred] (0,0)--(0,1) node[right]{$\colvec{0\\1}$};
  \node[align=center] at (1.3,-1.0) {$\mat{B}\mat{A}\,\vec{e}_1$\\ rotate \emph{then} stretch};
\end{scope}
\end{tikzpicture}
\caption{Non-commutativity you can see. The same starting vector $\vec{e}_1$ ends up in
different places depending on the order of the rotation and the stretch.}
\label{ch08:noncommute}
\end{figure}

\begin{trapbox}
Because $\mat{A}\mat{B} \ne \mat{B}\mat{A}$, ordinary algebra identities can fail. For
matrices,
\[
(\mat{A}+\mat{B})^2 = \mat{A}^2 + \mat{A}\mat{B} + \mat{B}\mat{A} + \mat{B}^2,
\]
\emph{not} $\mat{A}^2 + 2\mat{A}\mat{B} + \mat{B}^2$, because you may not merge
$\mat{A}\mat{B}$ with $\mat{B}\mat{A}$. Likewise $(\mat{A}\mat{B})^2 = \mat{A}\mat{B}\mat{A}\mat{B}$, which you cannot rearrange into $\mat{A}^2\mat{B}^2$. When you find
yourself ``swapping factors,'' stop and ask whether you earned it.
\end{trapbox}

\subsection{The rules that DO hold}

Multiplication is still well behaved in every way except commuting. For compatible
shapes and any scalar $c$:

\begin{itemize}
\item \textbf{Distributive (left):} $\ \mat{A}(\mat{B}+\mat{C}) = \mat{A}\mat{B} + \mat{A}\mat{C}$.
\item \textbf{Distributive (right):} $\ (\mat{A}+\mat{B})\mat{C} = \mat{A}\mat{C} + \mat{B}\mat{C}$.
\item \textbf{Scalars slide:} $\ c(\mat{A}\mat{B}) = (c\mat{A})\mat{B} = \mat{A}(c\mat{B})$.
\item \textbf{Associative:} $\ (\mat{A}\mat{B})\mat{C} = \mat{A}(\mat{B}\mat{C})$.
\end{itemize}

Each follows because the entries are built from sums and products of real numbers, where
these laws already hold; the matrix versions just package many scalar identities at once.
Keep both distributive laws separate in your mind --- with no commuting, ``left'' and
``right'' are genuinely different operations.

\section{The Identity: the Machine That Does Nothing}
\label{ch08:identity}

Every kind of multiplication has a ``$1$.'' For $n\times n$ matrices it is the
\emph{identity matrix} $\mat{I}$ --- ones on the diagonal, zeros elsewhere:

\[
\mat{I}_2 = \begin{bmatrix} 1 & 0 \\ 0 & 1 \end{bmatrix},
\qquad
\mat{I}_3 = \begin{bmatrix} 1 & 0 & 0 \\ 0 & 1 & 0 \\ 0 & 0 & 1 \end{bmatrix}.
\]

\begin{intuitionbox}
$\mat{I}$ is the machine that leaves every vector exactly where it is: $\mat{I}\vec{x} = \vec{x}$ for all $\vec{x}$. Its columns are the basis vectors $\vec{e}_1,\dots,\vec{e}_n$
sitting in their home positions. Geometrically it does nothing at all --- and that is
precisely why it acts as the ``$1$'' of matrix multiplication.
\end{intuitionbox}

For any matrix $\mat{A}$ of the right shape, $\mat{I}\mat{A} = \mat{A}$ and
$\mat{A}\mat{I} = \mat{A}$. Quick check via the column view: the $j$-th column of
$\mat{A}\mat{I}$ is $\mat{A}$ times the $j$-th column of $\mat{I}$, which is
$\mat{A}\vec{e}_j$ --- and $\mat{A}\vec{e}_j$ is just the $j$-th column of $\mat{A}$.
So $\mat{A}\mat{I}$ has the same columns as $\mat{A}$. Done.

\begin{remark}
$\mat{I}$ is the rare matrix that commutes with \emph{everything}: $\mat{A}\mat{I} = \mat{I}\mat{A} = \mat{A}$. So is any scalar multiple $c\mat{I}$ (a uniform scaling). Apart
from these, expecting two matrices to commute is the exception, not the rule.
\end{remark}

\section{The Inverse: Undoing a Machine}
\label{ch08:inverse}

We can now ask the chapter's second big question. The machine $\mat{A}$ sends $\vec{x}$
to $\mat{A}\vec{x}$. Is there a machine that sends it back?

\begin{definition}[Inverse matrix]
\label{ch08:inv-def}
A square matrix $\mat{A}$ is \emph{invertible} (or \emph{nonsingular}) if there is a
square matrix $\mat{A}\inv$ of the same size with
\[
\mat{A}\inv \mat{A} = \mat{I} \qquad\text{and}\qquad \mat{A}\,\mat{A}\inv = \mat{I}.
\]
If no such matrix exists, $\mat{A}$ is \emph{singular}.
\end{definition}

\begin{intuitionbox}
$\mat{A}\inv$ is the ``rewind'' machine. Run $\vec{x}$ forward through $\mat{A}$ to get
$\vec{b} = \mat{A}\vec{x}$; run $\vec{b}$ through $\mat{A}\inv$ and you are back at
$\vec{x}$, because $\mat{A}\inv\vec{b} = \mat{A}\inv\mat{A}\vec{x} = \mat{I}\vec{x} = \vec{x}$. The inverse exists exactly when the forward machine loses no information ---
when you could always tell, from the output, what the input must have been.
\end{intuitionbox}

This information-loss picture is the whole story of \emph{when} an inverse exists. If
$\mat{A}$ ever sends two different inputs to the same output, the rewind is impossible:
seeing the output, you cannot know which input it came from. So $\mat{A}$ is invertible
exactly when it is \emph{one-to-one} --- and for a square matrix, being one-to-one is the
same as hitting every target, which is the same as having independent columns, which is
the same as nonzero determinant. We collect these equivalences below.

\subsection{The $2\times2$ formula}

In two dimensions there is a clean closed form worth memorizing.

\begin{theorem}[Inverse of a $2\times2$]
\label{ch08:2x2inv}
Let $\mat{A} = \begin{bmatrix} a & b \\ c & d \end{bmatrix}$ and let
$\det\mat{A} = ad - bc$. If $\det\mat{A}\ne 0$, then
\[
\mat{A}\inv = \frac{1}{ad-bc}\begin{bmatrix} d & -b \\ -c & a \end{bmatrix}.
\]
If $ad - bc = 0$, then $\mat{A}$ has no inverse.
\end{theorem}

\begin{proof}
Multiply and see. Let $\mat{M} = \begin{bmatrix} d & -b \\ -c & a \end{bmatrix}$. Then
\[
\mat{A}\mat{M} = \begin{bmatrix} a & b \\ c & d \end{bmatrix}
\begin{bmatrix} d & -b \\ -c & a \end{bmatrix}
= \begin{bmatrix} ad - bc & -ab + ab \\ cd - dc & -bc + ad \end{bmatrix}
= (ad-bc)\,\mat{I}.
\]
So if $ad-bc \ne 0$ we may divide by it, and $\frac{1}{ad-bc}\mat{M}$ is the inverse.
If $ad - bc = 0$ then $\mat{A}\mat{M} = \mat{0}$ with $\mat{M}\ne\mat{0}$ unless
$\mat{A}=\mat{0}$; an invertible $\mat{A}$ could be cancelled to force $\mat{M}=\mat{0}$,
a contradiction --- so $\mat{A}$ cannot be invertible.
\end{proof}

\begin{labbox}
Invert $\mat{A} = \begin{bmatrix} 4 & 7 \\ 2 & 6 \end{bmatrix}$ by the formula. The
determinant is $4\cdot6 - 7\cdot2 = 24 - 14 = 10$. Swap the diagonal, negate the
off-diagonal, divide by $10$:
\[
\mat{A}\inv = \frac{1}{10}\begin{bmatrix} 6 & -7 \\ -2 & 4 \end{bmatrix}
= \begin{bmatrix} 0.6 & -0.7 \\ -0.2 & 0.4 \end{bmatrix}.
\]
Now \emph{check it} (always check):
\[
\mat{A}\mat{A}\inv = \frac{1}{10}\begin{bmatrix} 4 & 7 \\ 2 & 6 \end{bmatrix}
\begin{bmatrix} 6 & -7 \\ -2 & 4 \end{bmatrix}
= \frac{1}{10}\begin{bmatrix} 24-14 & -28+28 \\ 12-12 & -14+24 \end{bmatrix}
= \frac{1}{10}\begin{bmatrix} 10 & 0 \\ 0 & 10 \end{bmatrix} = \mat{I}.\ \checkmark
\]
\end{labbox}

\subsection{The geometry of singular: area collapse}

Why is $\det\mat{A} = ad - bc$ the make-or-break number? Because, as you will prove in
the determinant chapter, $\abs{\det\mat{A}}$ is the \emph{area} of the parallelogram
that $\mat{A}$ makes from the unit square --- the factor by which $\mat{A}$ scales area.

\begin{figure}[H]
\centering
\begin{tikzpicture}[scale=1.0,>=Stealth]
% invertible: unit square -> parallelogram with area
\begin{scope}[xshift=0cm]
  \draw[->,gray!45] (-0.3,0)--(3.4,0);
  \draw[->,gray!45] (0,-0.3)--(0,3.0);
  \fill[vecblue!18] (0,0)--(2,0.5)--(2.5,2)--(0.5,1.5)--cycle;
  \draw[->,thick,vecblue] (0,0)--(2,0.5);
  \draw[->,thick,vecblue] (0,0)--(0.5,1.5);
  \node at (1.55,0.95) {area $\ne 0$};
  \node[align=center] at (1.5,-1.0){invertible\\ $\det\ne0$};
\end{scope}
% singular: collapse to a line
\begin{scope}[xshift=6.2cm]
  \draw[->,gray!45] (-0.3,0)--(3.4,0);
  \draw[->,gray!45] (0,-0.3)--(0,3.0);
  \draw[very thick,vecred] (0,0)--(2.6,1.3);
  \draw[->,thick,vecred] (0,0)--(2,1.0);
  \draw[->,thick,vecred] (0,0)--(1.0,0.5);
  \node[vecred] at (2.0,2.0) {area $=0$};
  \node[align=center] at (1.5,-1.0){singular\\ $\det=0$};
\end{scope}
\end{tikzpicture}
\caption{Determinant as area. An invertible matrix turns the unit square into a
parallelogram with nonzero area. A singular matrix flattens it onto a line --- the
columns are parallel, area is zero, and the collapse cannot be undone.}
\label{ch08:areacollapse}
\end{figure}

\begin{intuitionbox}
If $\det\mat{A} = 0$, the matrix squashes the plane onto a line (or a point). A whole
direction's worth of vectors gets crushed to zero. Once an entire line of inputs lands on
a single output, no machine can pull them back apart --- the rewind would have to send
one point to many. \textbf{Zero determinant means lost information means no inverse.}
That is the geometry behind ``singular.''
\end{intuitionbox}

\subsection{The big equivalence (collect everything)}

For square matrices, a long list of statements all say the same thing. You have met some
already; the rest arrive in the chapters on subspaces, rank, and determinants. We gather
them so you can see the unity.

\begin{theorem}[The Invertible Matrix Theorem, first installment]
\label{ch08:imt}
For a square matrix $\mat{A}$ of size $n\times n$, the following are equivalent --- all
true together or all false together:
\begin{enumerate}
\item $\mat{A}$ is invertible.
\item $\mat{A}\vec{x} = \vec{b}$ has exactly one solution for every $\vec{b}$.
\item $\mat{A}\vec{x} = \vec{0}$ has only the trivial solution $\vec{x} = \vec{0}$.
\item The columns of $\mat{A}$ are linearly independent.
\item The columns of $\mat{A}$ span all of $\R^n$.
\item $\mat{A}$ row-reduces to the identity $\mat{I}$ (it has $n$ pivots).
\item $\det\mat{A} \ne 0$.
\end{enumerate}
\end{theorem}

We will not prove the whole chain here --- several links wait for later chapters --- but
the spine is the rewind picture. Invertibility (1) is exactly ``no information lost,''
which is exactly ``nothing nonzero gets crushed to zero'' (3), which is exactly
``independent columns'' (4). Everything else hangs from that spine.

\begin{strangbox}
Strang loves to organize the whole course around this equivalence; he often calls it
``the most important theorem in linear algebra'' and returns to it again and again,
adding a new equivalent statement each time a new concept appears (pivots, rank,
nullspace, eigenvalues, even singular values). \textbf{Before Strang:} read this list as
seven ways of saying ``the machine is reversible.'' \textbf{After Strang:} you will be
able to prove every arrow, and you will have added ``$0$ is not an eigenvalue'' and
``all singular values positive'' to the bottom of the list.
\end{strangbox}

\section{Computing the Inverse with Gauss--Jordan}
\label{ch08:gaussjordan}

The $2\times2$ formula does not scale --- the $3\times3$ closed form is already ugly, and
beyond that, hopeless. We need an algorithm. Happily it is just elimination, run with a
clever bookkeeping trick.

\subsection{The idea: solve $n$ systems at once}

What is the inverse, really? Its columns $\vec{x}_1,\dots,\vec{x}_n$ must satisfy

\[
\mat{A}\vec{x}_1 = \vec{e}_1,\quad
\mat{A}\vec{x}_2 = \vec{e}_2,\quad \dots,\quad
\mat{A}\vec{x}_n = \vec{e}_n,
\]

because then $\mat{A}[\,\vec{x}_1\ \cdots\ \vec{x}_n\,] = [\,\vec{e}_1\ \cdots\ \vec{e}_n\,]
= \mat{I}$, which says exactly that $[\,\vec{x}_1\ \cdots\ \vec{x}_n\,] = \mat{A}\inv$.
So finding the inverse is the same as solving $n$ linear systems that share the same
left-hand side $\mat{A}$. Since elimination on $\mat{A}$ is identical across all $n$
systems, we do it once on all the right-hand sides simultaneously.

\begin{focusbox}
\textbf{The Gauss--Jordan recipe.} Form the wide \emph{augmented matrix}
$[\,\mat{A}\mid\mat{I}\,]$. Row-reduce until the left block becomes $\mat{I}$. Whatever
the right block has become is $\mat{A}\inv$:
\[
[\,\mat{A}\mid \mat{I}\,] \ \xrightarrow{\text{row operations}}\ [\,\mat{I}\mid \mat{A}\inv\,].
\]
If the left block cannot be reduced to $\mat{I}$ (a zero row appears), then $\mat{A}$ is
singular and has no inverse.
\end{focusbox}

Why does the right block come out as $\mat{A}\inv$? Every row operation is itself
multiplication on the left by some elementary matrix $\mat{E}$. If the full sequence is
$\mat{E} = \mat{E}_k\cdots\mat{E}_2\mat{E}_1$, then reducing the left block to $\mat{I}$
means $\mat{E}\mat{A} = \mat{I}$, so $\mat{E} = \mat{A}\inv$. But the \emph{same}
$\mat{E}$ acts on the right block, turning $\mat{I}$ into $\mat{E}\mat{I} = \mat{E} =
\mat{A}\inv$. The augmented matrix is a machine for applying $\mat{A}\inv$ to both
blocks at once; we just watch what it does to the identity.

\subsection{A full $3\times3$ worked example}

Let us invert

\[
\mat{A} = \begin{bmatrix} 2 & 1 & 1 \\ 1 & 3 & 2 \\ 1 & 0 & 0 \end{bmatrix}.
\]

Start by augmenting with $\mat{I}$:

\[
\left[\begin{array}{ccc|ccc}
2 & 1 & 1 & 1 & 0 & 0 \\
1 & 3 & 2 & 0 & 1 & 0 \\
1 & 0 & 0 & 0 & 0 & 1
\end{array}\right].
\]

\textbf{Step 1.} The bottom row is the simplest, so swap rows $1$ and $3$ to put a clean
pivot $1$ in the top-left ($R_1 \leftrightarrow R_3$):

\[
\left[\begin{array}{ccc|ccc}
1 & 0 & 0 & 0 & 0 & 1 \\
1 & 3 & 2 & 0 & 1 & 0 \\
2 & 1 & 1 & 1 & 0 & 0
\end{array}\right].
\]

\textbf{Step 2.} Clear the rest of column $1$. Replace $R_2 \to R_2 - R_1$ and
$R_3 \to R_3 - 2R_1$:

\[
\left[\begin{array}{ccc|ccc}
1 & 0 & 0 & 0 & 0 & 1 \\
0 & 3 & 2 & 0 & 1 & -1 \\
0 & 1 & 1 & 1 & 0 & -2
\end{array}\right].
\]

\textbf{Step 3.} We want a clean pivot in column $2$. Row $3$ has a $1$ there, so swap
$R_2 \leftrightarrow R_3$:

\[
\left[\begin{array}{ccc|ccc}
1 & 0 & 0 & 0 & 0 & 1 \\
0 & 1 & 1 & 1 & 0 & -2 \\
0 & 3 & 2 & 0 & 1 & -1
\end{array}\right].
\]

\textbf{Step 4.} Clear below the second pivot: $R_3 \to R_3 - 3R_2$:

\[
\left[\begin{array}{ccc|ccc}
1 & 0 & 0 & 0 & 0 & 1 \\
0 & 1 & 1 & 1 & 0 & -2 \\
0 & 0 & -1 & -3 & 1 & 5
\end{array}\right].
\]

\textbf{Step 5.} Make the third pivot $+1$: $R_3 \to -R_3$:

\[
\left[\begin{array}{ccc|ccc}
1 & 0 & 0 & 0 & 0 & 1 \\
0 & 1 & 1 & 1 & 0 & -2 \\
0 & 0 & 1 & 3 & -1 & -5
\end{array}\right].
\]

\textbf{Step 6.} Now eliminate \emph{upward} to clear above the third pivot. Only $R_2$
has an entry there: $R_2 \to R_2 - R_3$:

\[
\left[\begin{array}{ccc|ccc}
1 & 0 & 0 & 0 & 0 & 1 \\
0 & 1 & 0 & -2 & 1 & 3 \\
0 & 0 & 1 & 3 & -1 & -5
\end{array}\right].
\]

The left block is now $\mat{I}$ (column $1$ was already clean above the pivot). Read off
the right block:

\[
\mat{A}\inv = \begin{bmatrix} 0 & 0 & 1 \\ -2 & 1 & 3 \\ 3 & -1 & -5 \end{bmatrix}.
\]

\textbf{Always verify.} Multiply $\mat{A}\mat{A}\inv$ and demand $\mat{I}$:

\[
\begin{bmatrix} 2 & 1 & 1 \\ 1 & 3 & 2 \\ 1 & 0 & 0 \end{bmatrix}
\begin{bmatrix} 0 & 0 & 1 \\ -2 & 1 & 3 \\ 3 & -1 & -5 \end{bmatrix}.
\]
Row $1$: $(2\cdot0 + 1\cdot(-2) + 1\cdot3,\ \ 2\cdot0 + 1\cdot1 + 1\cdot(-1),\ \
2\cdot1 + 1\cdot3 + 1\cdot(-5)) = (1,0,0)$. Row $2$:
$(0 - 6 + 6,\ 0 + 3 - 2,\ 1 + 9 - 10) = (0,1,0)$. Row $3$:
$(0 + 0 + 0,\ 0,\ 1 + 0 + 0) = (0,0,1)$. Together that is exactly $\mat{I}$. $\checkmark$

\begin{trapbox}
Two classic Gauss--Jordan slips. First, \textbf{do the row operation to the whole row} ---
both blocks, every time. It is fatally easy to update the left and forget the right.
Second, watch your \textbf{direction}: ordinary elimination only clears \emph{below} each
pivot (giving an upper-triangular $\mat{U}$); for the inverse you must also clear
\emph{above}, all the way to a clean $\mat{I}$. ``Gauss'' goes down; ``Jordan'' adds the
trip back up.
\end{trapbox}

\begin{labbox}
Your turn, by hand. Invert $\mat{A} = \begin{bmatrix} 1 & 2 & 0 \\ 0 & 1 & 1 \\ 0 & 0 & 1 \end{bmatrix}$ with $[\,\mat{A}\mid\mat{I}\,]$. It is upper triangular, so you only ever
clear upward; no row swaps needed. (Answer:
$\mat{A}\inv = \begin{bmatrix} 1 & -2 & 2 \\ 0 & 1 & -1 \\ 0 & 0 & 1 \end{bmatrix}$ ---
check it.) Notice how the $-2$ and the $+2$ track the two elimination steps; the inverse
remembers exactly what you undid.
\end{labbox}

\subsection{What singular looks like in the algorithm}

If $\mat{A}$ is singular, elimination cannot manufacture a missing pivot. Somewhere a
row of the left block becomes all zeros. The moment that happens, stop --- there is no
inverse. Try it on $\mat{A} = \begin{bmatrix} 1 & 2 \\ 2 & 4 \end{bmatrix}$ (note row
$2$ is twice row $1$, so $\det = 4 - 4 = 0$):

\[
\left[\begin{array}{cc|cc}
1 & 2 & 1 & 0 \\ 2 & 4 & 0 & 1
\end{array}\right]
\xrightarrow{R_2 \to R_2 - 2R_1}
\left[\begin{array}{cc|cc}
1 & 2 & 1 & 0 \\ 0 & 0 & -2 & 1
\end{array}\right].
\]

The left block has a zero row; it can never become $\mat{I}$. The algorithm itself
\emph{detects} singularity --- you do not need to compute a determinant first. The
collapse you saw geometrically (the two columns $\colvec{1\\2}$ and $\colvec{2\\4}$ are
parallel, so the unit square flattens onto a line) shows up as that fatal zero row.

\section{The Inverse of a Product: Socks and Shoes}
\label{ch08:product-inverse}

We close with one identity you will use constantly, and it has the most memorable mnemonic
in the subject.

\begin{theorem}[Inverse of a product]
\label{ch08:abinv}
If $\mat{A}$ and $\mat{B}$ are both invertible (and the same size), then $\mat{A}\mat{B}$
is invertible and
\[
(\mat{A}\mat{B})\inv = \mat{B}\inv \mat{A}\inv.
\]
The order reverses.
\end{theorem}

\begin{proof}
Just multiply the proposed inverse against $\mat{A}\mat{B}$ and watch it collapse, using
associativity to regroup:
\[
(\mat{A}\mat{B})(\mat{B}\inv\mat{A}\inv)
= \mat{A}(\mat{B}\mat{B}\inv)\mat{A}\inv
= \mat{A}\,\mat{I}\,\mat{A}\inv
= \mat{A}\mat{A}\inv = \mat{I}.
\]
The same cancellation from the other side gives
$(\mat{B}\inv\mat{A}\inv)(\mat{A}\mat{B}) = \mat{I}$. By
Definition~\ref{ch08:inv-def}, $\mat{B}\inv\mat{A}\inv$ is the inverse of $\mat{A}\mat{B}$.
\end{proof}

\begin{intuitionbox}
\textbf{Socks and shoes.} In the morning you put on socks ($\mat{B}$), then shoes
($\mat{A}$). To undo it at night you cannot take off the socks first --- the shoes are in
the way. You reverse the order: shoes off ($\mat{A}\inv$), then socks off ($\mat{B}\inv$).
To undo a composition, undo the \emph{last} thing first. That is exactly
$(\mat{A}\mat{B})\inv = \mat{B}\inv\mat{A}\inv$.
\end{intuitionbox}

\begin{figure}[H]
\centering
\begin{tikzpicture}[node distance=8mm and 16mm,>=Stealth,font=\small]
\node (x) {$\vec{x}$};
\node[right=of x] (bx) {$\mat{B}\vec{x}$};
\node[right=of bx] (abx) {$\mat{A}\mat{B}\vec{x}$};
\draw[->,thick] (x)--(bx) node[midway,above]{$\mat{B}$};
\draw[->,thick] (bx)--(abx) node[midway,above]{$\mat{A}$};
\draw[->,thick,vecred,bend left=38] (abx) to node[midway,below]{$\mat{A}\inv$} (bx);
\draw[->,thick,vecred,bend left=38] (bx) to node[midway,below]{$\mat{B}\inv$} (x);
\end{tikzpicture}
\caption{Forward (black) you apply $\mat{B}$ then $\mat{A}$. To rewind (red) you must
peel them off in reverse: $\mat{A}\inv$ first, then $\mat{B}\inv$. Hence
$(\mat{A}\mat{B})\inv = \mat{B}\inv\mat{A}\inv$.}
\label{ch08:socksshoes}
\end{figure}

The pattern extends to any chain: $(\mat{A}\mat{B}\mat{C})\inv = \mat{C}\inv\mat{B}\inv\mat{A}\inv$, reverse the whole order. The same reversal governs transposes,
$(\mat{A}\mat{B})\T = \mat{B}\T\mat{A}\T$, which you will meet in the next chapter ---
two different operations, one shared ``order reverses'' rhythm.

\begin{connectionbox}
\textbf{Groups.} The invertible $n\times n$ matrices, with multiplication, form a
\emph{group} --- the \emph{general linear group} $GL_n(\R)$. The four axioms are exactly
what this chapter established: multiplication is associative (Theorem~\ref{ch08:assoc}),
there is an identity $\mat{I}$, every element has an inverse, and the product of two
invertibles is invertible with $(\mat{A}\mat{B})\inv = \mat{B}\inv\mat{A}\inv$. The
order-reversing inverse rule is a \emph{general} feature of groups, not a quirk of
matrices --- it holds for symmetries of a cube, permutations of a deck, and the
rotations of space. When physicists speak of symmetry groups and Lie groups underlying
the laws of nature, $GL_n$ and its subgroups (rotations $SO_n$, volume-preserving maps
$SL_n$) are the founding examples. You have just verified the group axioms by hand.
\end{connectionbox}

\begin{remark}[A useful shortcut you have earned]
To \emph{prove} a square matrix $\mat{C}$ is invertible with inverse $\mat{D}$, you do
not have to check both $\mat{C}\mat{D}=\mat{I}$ and $\mat{D}\mat{C}=\mat{I}$. For square
matrices, either one alone implies the other (a fact that follows from the equivalence in
Theorem~\ref{ch08:imt}). That is why the one-sided check in the proof above was enough.
Be careful, though: this rescue is for \emph{square} matrices only. Non-square matrices
can have a one-sided inverse on the left or the right but never both at once.
\end{remark}

\section*{Looking Back, Looking Ahead}

Two questions, two answers. \emph{What is two machines in a row?} A single matrix
$\mat{A}\mat{B}$, best understood through the four views --- entries, columns, rows, and
rank-one layers --- and obeying every algebra rule except commuting. \emph{Can we undo a
machine?} Only when it loses no information; then the rewind is $\mat{A}\inv$, computed
by Gauss--Jordan, and products rewind in reverse order.

\begin{focusbox}
\textbf{Cost and intuition, side by side.} The dot-product view tells you the
\emph{price}: an $n\times n$ product is about $n^3$ multiply--adds, and inverting an
$n\times n$ matrix by Gauss--Jordan is the same $n^3$ order of work (elimination on $n$
right-hand sides at once). The other three views tell you the \emph{meaning}: columns of
$\mat{A}\mat{B}$ live in $\col(\mat{A})$, rows live in $\row(\mat{B})$, and the whole
product is a sum of rank-one layers. Price comes from View~1; understanding --- and every
shortcut that beats the naive $n^3$ --- comes from Views 2--4.
\end{focusbox}

These tools turn the next several chapters into a story of \emph{factoring} machines into
simpler ones. $\mat{A} = \mat{L}\mat{U}$ records elimination itself as a product.
$\mat{A} = \mat{Q}\mat{R}$ separates rotation from stretch. And the crown jewel,
$\mat{A} = \mat{U}\mat{\Sigma}\mat{V}\T$, says \emph{every} matrix is rotate, then
stretch along axes, then rotate again --- three simple machines whose product is any
linear map at all. All of it is built on the multiplication and inverses you
now own.

\section{Exercises}
\label{ch08:exercises}

\begin{exercise}[Warm-up by hand]
Compute $\mat{A}\mat{B}$ and $\mat{B}\mat{A}$ for
$\mat{A} = \begin{bmatrix} 1 & 2 \\ 0 & 1 \end{bmatrix}$,
$\mat{B} = \begin{bmatrix} 1 & 0 \\ 3 & 1 \end{bmatrix}$. Confirm they are different, and
say in one sentence why two shears in different orders give different results.
\end{exercise}

\begin{exercise}[All four views]
For $\mat{A} = \begin{bmatrix} 2 & 0 \\ 1 & 3 \end{bmatrix}$ and
$\mat{B} = \begin{bmatrix} 1 & 4 \\ 2 & 0 \end{bmatrix}$, compute $\mat{A}\mat{B}$
\emph{four ways}: (a) entry by entry (dot products), (b) as combinations of $\mat{A}$'s
columns, (c) as combinations of $\mat{B}$'s rows, (d) as a sum of two outer products.
Verify all four agree.
\end{exercise}

\begin{exercise}[Shapes]
$\mat{A}$ is $3\times 2$, $\mat{B}$ is $2\times 5$, $\mat{C}$ is $5\times 3$. Which of
$\mat{A}\mat{B}$, $\mat{B}\mat{C}$, $\mat{C}\mat{A}$, $\mat{A}\mat{B}\mat{C}$,
$\mat{C}\mat{A}\mat{B}$ are defined, and what is each one's shape? Which products
$\mat{X}\mat{Y}$ from this list are square?
\end{exercise}

\begin{exercise}[$2\times2$ inverse]
Use the formula to invert $\mat{A} = \begin{bmatrix} 3 & 5 \\ 1 & 2 \end{bmatrix}$, then
check $\mat{A}\mat{A}\inv = \mat{I}$. What is $\det\mat{A}$, and how does it appear in the
answer?
\end{exercise}

\begin{exercise}[Gauss--Jordan]
Invert $\mat{A} = \begin{bmatrix} 1 & 1 & 1 \\ 0 & 2 & 1 \\ 1 & 0 & 1 \end{bmatrix}$ by
augmenting with $\mat{I}$ and row-reducing to $[\,\mat{I}\mid\mat{A}\inv\,]$. Show every
step, and verify your answer.
\end{exercise}

\begin{exercise}[Detecting singular]
Run Gauss--Jordan on $\mat{A} = \begin{bmatrix} 1 & 2 & 3 \\ 2 & 4 & 6 \\ 0 & 1 & 1 \end{bmatrix}$ far enough to show it is singular. Identify the row that forces the
collapse, and describe geometrically what $\mat{A}$ does to the unit cube.
\end{exercise}

\begin{exercise}[Socks and shoes, with numbers]
Let $\mat{A} = \begin{bmatrix} 2 & 0 \\ 0 & 1 \end{bmatrix}$ and
$\mat{B} = \begin{bmatrix} 1 & 1 \\ 0 & 1 \end{bmatrix}$. Compute $\mat{A}\inv$,
$\mat{B}\inv$, and $\mat{A}\mat{B}$. Verify directly that $(\mat{A}\mat{B})\inv =
\mat{B}\inv\mat{A}\inv$, and check that $\mat{A}\inv\mat{B}\inv$ gives the \emph{wrong}
answer (it is not the inverse of $\mat{A}\mat{B}$).
\end{exercise}

\begin{exercise}[Rank cannot grow]
Using the column view, explain why every column of $\mat{A}\mat{B}$ is a linear
combination of the columns of $\mat{A}$. Conclude that $\rank(\mat{A}\mat{B}) \le
\rank(\mat{A})$. Then give a $2\times2$ example where $\mat{A}\ne\mat{0}$, $\mat{B}\ne
\mat{0}$, yet $\mat{A}\mat{B} = \mat{0}$ --- so rank can \emph{drop} all the way to zero.
\end{exercise}

\begin{exercise}[Inverse of a diagonal matrix]
Let $\mat{D} = \diag(d_1,\dots,d_n)$. Show that $\mat{D}$ is invertible if and only if
every $d_i \ne 0$, and that then $\mat{D}\inv = \diag(1/d_1,\dots,1/d_n)$. What goes
wrong, geometrically and in Gauss--Jordan, if some $d_i = 0$?
\end{exercise}

\begin{exercise}[Proof: inverses are unique]
Suppose $\mat{B}$ and $\mat{C}$ both satisfy $\mat{B}\mat{A} = \mat{I}$ and
$\mat{A}\mat{C} = \mat{I}$. Prove $\mat{B} = \mat{C}$. (Hint: simplify the product
$\mat{B}\mat{A}\mat{C}$ in two ways, using associativity.) Conclude that a matrix has at
most one inverse.
\end{exercise}

\begin{exercise}[Challenge: a $2\times2$ identity]
For an invertible $2\times2$ matrix $\mat{A} = \begin{bmatrix} a & b \\ c & d \end{bmatrix}$, prove the surprisingly clean fact that
\[
\mat{A}\inv = \frac{1}{\det\mat{A}}\bigl((\tr\mat{A})\,\mat{I} - \mat{A}\bigr),
\]
where $\tr\mat{A} = a + d$. (This is the $2\times2$ case of the Cayley--Hamilton
theorem, which you will meet again after eigenvalues. Just expand both sides.)
\end{exercise}

\begin{exercise}[Challenge: commuting with everything]
Prove that if a square matrix $\mat{M}$ satisfies $\mat{M}\mat{A} = \mat{A}\mat{M}$ for
\emph{every} matrix $\mat{A}$ of the same size, then $\mat{M} = c\mat{I}$ for some scalar
$c$. (Hint: try $\mat{A} = \vec{e}_i\vec{e}_j\T$, the matrix with a single $1$ in
position $(i,j)$, and see what commuting forces on the entries of $\mat{M}$.)
\end{exercise}
